\originalchapter{推理、搜索与可验证反馈}\label{ch:reason}
\originalsection{行为能力与内部过程}
“推理”需要先定义观测对象. 本书在行为层面指能够按给定规则组合多步信息, 得到可验证的结论, 并在新的组合和有控制的扰动上保持正确. 这一定义不要求输出长段解释, 也不把顺畅语言当作成功标准. 对形式证明可检查每一步推导, 对程序可运行测试, 对开放问题则需要更难的证据收集与判断.

思维链提示在部分算术、常识和符号任务中改善了大模型的报告成绩\cite{cot}. 中间 token 可以提供额外计算, 保存子结果, 并改变后续条件分布. 但可见解释未必忠实反映全部内部原因, 也可能是在得到答案后生成合理化文本. 对“理由是否忠实”的检验应改变某个真正影响答案的因素, 看解释是否追踪它, 并与只看输出的基线比较.

预测目标不禁止算法式计算. 一个模型可以把状态保存在隐向量中, 按多层变换完成某种算法. 但有限 benchmark 的正确答案, 仍不能唯一确定是哪一种机制: 记忆、相似例子匹配、启发式或可靠算法都可能产生相同输出. 组合留出、长度外推与反事实干预有助于辨别, 每种检验也都有覆盖限制.

\originalsection{采样、验证与选择}
若每个候选独立地以概率 $p$ 正确, $N$ 个候选中至少一个正确的概率为 $1-(1-p)^N$. 这个式子是生成覆盖率, 不等于最终选中答案的正确率. 若验证器从 $N$ 个候选中无法可靠挑选, 多生成可能只得到更多高分错误. 候选之间强相关时, 独立式也失效.

\begin{definition}
本书把 pass@$N$ 理解为给定 $N$ 次生成中存在正确答案的比例, 把 selected accuracy 理解为实际选择器提交的那个答案的正确率. 两者必须分别评价, 同时报告候选预算、采样规则和验证器的访问能力.
\end{definition}
对精确可检验任务, 验证器的设计可能比生成器更容易. 例如解整数方程 $a^2+b^2=65$, 检查一个候选只需两次平方与求和. 在更复杂证明中, 验证器也可能存在边界情况遗漏、数值误差或规格不全. 一份程序通过有限单元测试不等于对所有输入正确; 形式验证又需要证明规格确实表达了目标.

本项目 \texttt{search\_verifier} 从 $[-9,9]^2$ 中生成固定的 600 个候选, 对预算 $1,4,16,64,256$ 各作 200 次池内不放回采样, 检查是否存在满足方程的候选. 这里没有 LLM, 只演示“生成覆盖加精确反馈”的机制. 各次试验共享候选池, 不能冒称独立模型试验. 若池内共有 $K$ 个正确候选、总量 $M$, 精确无正确概率为 $\binom{M-K}{N}/\binom MN$, 而不是有放回的 $(1-K/M)^N$. 组合数不可能时无正确概率取零.

\originalsection{奖励与策略梯度}
把语言模型视为条件策略 $\pi_\theta(y\mid x)$, 序列奖励为 $r(x,y)$. 对固定有限候选集、参数可微、支持不随参数改变的情形,
\begin{align*}
\nabla_\theta\E_{y\sim\pi_\theta}[r(y)]
 &=\sum_y r(y)\nabla_\theta\pi_\theta(y)\\
 &=\E_{y\sim\pi_\theta}[r(y)\nabla_\theta\log\pi_\theta(y)].
\end{align*}
若序列自回归生成, $\log\pi_\theta(y\mid x)$ 是各 token 对数概率之和, 于是序列奖励影响整条轨迹的更新. 可减去不依赖 $y$ 的基线 $b(x)$, 因为 $\E[\nabla\log\pi]=\nabla\sum_y\pi(y)=0$. 基线可以降低方差, 不改变理想期望梯度.

奖励正确答案会让高奖励轨迹更可能出现. 它不会自动保证每个中间步骤正确; 若奖励器有漏洞, 模型可能提高漏洞利用率而不是任务能力. 数学、程序等具有较明确验证器的任务适合这类研究, 开放世界判断则更难. DeepSeek-R1 报告基于已有预训练模型, 通过可验证任务奖励和多阶段训练改善推理表现\cite{r1}. R1-Zero 缺少 RL 前推理监督微调, 不等于没有预训练知识、人类设计或大量计算. R1 的完整流程又不同于 R1-Zero.

2024 年关于测试时计算分配的研究表明, 预算应按问题难度与选择策略分配, 不能只比较“长回答”与“短回答”\cite{testcompute}. 如果模型在简单题上已可靠, 多次采样可能浪费资源; 如果错误模式高度相关, 延长搜索也不一定有收益. 推理 token 数、工具成本、墙钟时间和基础模型训练成本需要分别记账.

\originalsection{练习与解答}
\begin{exercise}
候选独立正确率为 $0.2$, 生成 4 次的至少一次正确概率是多少? 若验证器始终选择第一个, 最终正确率是多少?
\end{exercise}
\begin{solution}
覆盖率为 $1-0.8^4=0.5904$. 始终选第一个的正确率仍为 $0.2$. 因此 59.04\% 的 pass@4 并不意味着用户最终能得到同等正确率.
\end{solution}
\begin{exercise}
如果验证器只检查“程序成功退出”, 为什么搜索可能得到不解决目标的程序? 给出一种补充机制.
\end{exercise}
\begin{solution}
恒返回固定值也能成功退出, 与目标无关. 应把输出规格、边界条件和性质写成独立检查, 并设置隐藏测试或形式化验证. 验证器不得被候选程序修改. 补充检查也需审查它自身是否正确, 不能把“有测试”当成任意程序的正确性证明.
\end{solution}
\label{ch:end:10}
